\documentclass[10pt]{article}
\usepackage[utf8]{inputenc}
\usepackage[T1]{fontenc}
\usepackage{lmodern}
\usepackage[a4paper,margin=0.7in,landscape]{geometry}
\usepackage{longtable}
\usepackage{booktabs}
\usepackage{array}
\usepackage{microtype}
\usepackage[colorlinks=true,linkcolor=black,urlcolor=blue!55!black]{hyperref}
\setlength{\parindent}{0pt}
\renewcommand{\arraystretch}{1.15}

\begin{document}

\begin{center}
{\Large\bfseries Supplement 1. TRIPOD+AI reporting checklist}\\[4pt]
{\large AtrialSpectralBench: incremental predictive value and
perturbation-theoretic validity of a spectral fragility index for atrial reentry}\\[2pt]
Mounish Mavuduru
\end{center}

\vspace{0.6em}

\textbf{Scope of this checklist.} TRIPOD+AI adherence is claimed for
\emph{one} analysis in this manuscript: the pre-registered clinical endpoint
predicting two-year post-ablation arrhythmia recurrence in the 82 University of
Washington patients (Results, ``The clinical endpoint, and the false positive its
own protocol caught''). That is the only analysis scored against a real patient
outcome. Every other endpoint in the paper is in silico, with the outcome an
electrophysiology simulator's reentry-inducibility verdict rather than a patient
outcome; those analyses fall outside TRIPOD+AI's remit and no adherence is
claimed for them.

\textbf{How to read the status column.} This manuscript was written as a
pre-registered methodological evaluation ending in a reported null, not as a
clinical prediction-model development study, and the checklist reflects that
honestly rather than being padded. ``Partially'' means the manuscript addresses
the substance of the item but not in the form TRIPOD+AI specifies, or addresses
it for the analysis but not for the cohort's clinical characteristics.
``Not reported'' is recorded where the manuscript genuinely does not address the
item --- most often because no patient-level demographic or clinical data were
released with either public deposit. Items marked \emph{n/a} do not apply to a
single-cohort study with no model updating and no external validation set.

\textbf{Reference.} Collins GS, Moons KGM, Dhiman P, et al. TRIPOD+AI statement:
updated guidance for reporting clinical prediction models that use regression or
machine learning methods. \emph{BMJ} 2024;385:e078378. TRIPOD+AI supersedes the
2015 TRIPOD statement.

\vspace{0.8em}

\begin{longtable}{@{}p{0.9cm} p{6.2cm} p{1.9cm} p{12.4cm}@{}}
\toprule
\textbf{Item} & \textbf{TRIPOD+AI checklist item} & \textbf{Status} &
\textbf{Where reported, and comment} \\
\endfirsthead
\toprule
\textbf{Item} & \textbf{TRIPOD+AI checklist item} & \textbf{Status} &
\textbf{Where reported, and comment} \\
\endhead
\midrule
\multicolumn{4}{@{}l}{\textbf{Title}}\\
\midrule
1 & Identify the study as developing or evaluating the performance of a multivariable prediction model, the target population, and the outcome to be predicted & Partially & \textbf{Title page} --- The title identifies the work as ''a pre-registered evaluation'' of the ''incremental predictive value'' of a spectral fragility index for atrial reentry on 182 patient-derived left-atrial meshes, but it does not use prediction-model language, and neither the TRIPOD-scoped clinical target population (82 University of Washington atrial fibrillation ablation patients) nor the clinical outcome predicted (2-year post-ablation arrhythmia recurrence) appears in the title, which names only the in-silico reentry endpoint. \\[2pt]
\midrule
\multicolumn{4}{@{}l}{\textbf{Abstract}}\\
\midrule
2 & See TRIPOD+AI for Abstracts checklist & Partially & \textbf{Abstract, clinical sentences at; scope note} --- The unstructured abstract does report the pre-registered clinical endpoint on 2-year recurrence in the 82 UW patients with its result (dAUC +0.197, p=0.0061, endpoint not met, increment classified as an artefact of an anti-predictive AUC 0.250 baseline reproduced by nine noise columns at empirical p=0.333), and a standalone note states that this is the only analysis scored against real patient outcomes, but it omits most TRIPOD-for-Abstracts elements for that endpoint: no structured headings, no study design or data source, no number of outcome events, no predictors, no model type, and no statement of the validation procedure. \\[2pt]
\midrule
\multicolumn{4}{@{}l}{\textbf{Introduction}}\\
\midrule
3a & Explain the healthcare context (including whether diagnostic or prognostic) and rationale for developing or evaluating the prediction model, including references to existing models & Partially & \textbf{Introduction; Methods; clinical endpoint} --- The Introduction gives a clear mechanistic rationale for testing the spectral fragility index as a reentry-fragility biomarker and pre-registers the null, and the competitor feature set it must beat is named in Methods, but the manuscript never frames the clinical endpoint as prognostic, never sets out the healthcare context or clinical need for predicting post-ablation recurrence, and cites no existing clinical prediction models for atrial fibrillation recurrence (the AF-recurrence literature cited atis used only to benchmark simulated inducibility rates). \\[2pt]
3b & Describe the target population and the intended purpose of the prediction model in the context of the care pathway, including its intended users (e.g., healthcare professionals, patients, public) & Partially & \textbf{Methods, Cohort flow and Substrate; clinical endpoint} --- The analysed population is described as a dataset (82 distinct atrial fibrillation patients from the UW/Boyle LGE-MRI deposit, pre-ablation meshes only, with 2-year post-ablation outcomes supplied by the cohort authors), but the manuscript states no intended purpose for the index in the care pathway, no clinical decision it would inform, and no intended users; the work is presented throughout as a methodological evaluation ending in a negative result rather than as a model proposed for clinical use. \\[2pt]
3c & Describe any known health inequalities between sociodemographic groups Section/Topic Item Checklist item & Partially & \textbf{Limitations, item on sociodemographic characterisation} --- The manuscript now states explicitly that neither public deposit released any age, sex, race, ethnicity, comorbidity or socioeconomic variable, so the cohort cannot be characterised demographically, performance cannot be reported within sociodemographic subgroups, and representativeness relative to a target ablation population is unknown; it does not additionally review known health inequalities in ablation outcomes, so the item is met only in part. \\[2pt]
4 & Specify the study objectives, including whether the study describes the development or validation of a prediction model (or both) & Partially & \textbf{Introduction; Methods; clinical endpoint; non-independent split noted at} --- The objective for the clinical endpoint is stated precisely and was pre-registered (test whether a 9-column SFI block adds incremental discrimination over a 6-feature competitor baseline for 2-year recurrence, with a minimum detectable effect fixed at dAUC 0.10-0.15 and a fixed-sequence stopping rule), and the manuscript makes clear that no independent validation cohort exists (the provider's 67/15 split is explicitly labelled non-independent at, but the objectives are never framed in TRIPOD's terms of model development versus validation, and the study is not labelled as model development with internal validation. \\[2pt]
\midrule
\multicolumn{4}{@{}l}{\textbf{Methods}}\\
\midrule
5a & Describe the sources of data separately for the development and evaluation datasets (e.g., randomised trial, cohort, routine care or registry data), the rationale for using these data, and representat & Partially & \textbf{Methods, Cohort flow and Substrate; clinical endpoint; Reproducibility; Limitations} --- The single data source for the clinical endpoint is fully identified (UW/Boyle LGE-MRI cohort, Dryad 10.5061/dryad.kkwh70sg0, 82 distinct patients, 82 pre-ablation meshes with the 82 post-ablation meshes excluded because they encode the delivered treatment, outcomes supplied by the cohort authors), and a key data-quality limitation is disclosed (fibrosis is categorical because the intensity ratios are IRB-withheld, which degrades the baseline;, but no separate development and evaluation datasets exist, and the manuscript does not address whether these 82 patients are representative of the target population, since no patient-level clinical or demographic characteristics are reported at all. \\[2pt]
5b & Specify the dates of the collected participant data, including start and end of participant accrual; and, if applicable, end of follow-up & \emph{Not reported} & No calendar dates are given anywhere for the UW cohort: the manuscript specifies the follow-up rule and its length (recurrence as at least 30 s of documented atrial fibrillation or flutter after a 90-day blanking period, with every patient followed the full two years and no censoring, but reports no accrual start or end dates, no imaging dates, and no end-of-follow-up date. \\[2pt]
6a & Specify key elements of the study setting (e.g., primary care, secondary care, general population) including the number and location of centres & Partially & \textbf{Methods, ''Reporting guideline'' para and ''Substrate'' para; Results Sec. 5.2 sec:clinical; Reproducibility} --- The manuscript identifies the clinical cohort as the 82 University of Washington (UW/Boyle) LGE-MRI atrial fibrillation patients obtained from a single public Dryad deposit (10.5061/dryad.kkwh70sg0) with outcomes supplied by that cohort's authors, but it does not state the level of care, the number or location of recruiting centres, or the recruitment/imaging dates, all of which are left to the source publication. \\[2pt]
6b & Describe the eligibility criteria for study participants & Partially & \textbf{Methods, ''Cohort flow'' para; Methods ''Substrate'' para} --- Analysis-level eligibility is fully explicit -- all 82 UW patients in the deposit are analysed using their pre-ablation meshes only, ''No subject is excluded for any other reason,'' and there is no attrition -- but the underlying clinical eligibility criteria (AF type, ablation indication, imaging quality, age) are not restated and are deferred to the source cohort's citation. \\[2pt]
6c & Give details of any treatments received, and how they were handled during model development or evaluation, if relevant & Partially & \textbf{Methods, ''Cohort flow'' para; Results sec:clinical} --- All 82 patients underwent catheter ablation and the treatment is handled by design -- the primary analysis uses pre-ablation geometry only ''since post-ablation geometry encodes the delivered treatment,'' the outcome is defined post-ablation after a 90-day blanking period, and the pre-specified ablation-induced dSFI test was not performed under the fixed-sequence stopping rule -- but the ablation strategy, lesion set, energy source, and any concomitant antiarrhythmic drug therapy during follow-up are not described. \\[2pt]
7 & Describe any data pre-processing and quality checking, including whether this was similar across relevant sociodemographic groups & Partially & \textbf{Methods, ''The edge weight'' and ''Substrate''; Results Sec. 5.6 sec:substrateincl. 900, 925)} --- Pre-processing of the imaging-derived substrate is described in full (edge-weight construction, derivation of UW fibrosis from per-element tissue tags rather than the IRB-withheld intensity ratios, a PCA chart substituted for absent atrial coordinates, and vertex-clustering coarsening to \textasciitilde{}2000-node graphs) and quality checking is unusually detailed (a provenance audit that corrects the tag-164 cap misassignment and the constant-vector fibre placeholder, an Euler-characteristic census, and a loader that now refuses any mesh carrying an undeclared dead tag), but the manuscript reports no sociodemographic variables for the 82 patients and therefore never examines whether pre-processing behaved similarly across sociodemographic groups. \\[2pt]
8a & Clearly define the outcome that is being predicted and the time horizon, including how and when assessed, the rationale for choosing this outcome, and whether the method of outcome assessment is consi & Partially & \textbf{Results Sec. 5.2 sec:clinical; Methods ''Labels'' para; Reproducibility} --- The clinical outcome is defined precisely -- at least 30 s of documented atrial fibrillation or flutter after a 90-day blanking period, over a fixed two-year post-ablation horizon with every one of the 82 patients followed the full two years (no censoring), with the 48 events further split into 35 fibrillation and 13 flutter, and explicitly distinguished throughout from the simulator inducibility verdicts used elsewhere -- but the modality and schedule of rhythm monitoring used to document recurrence, the rationale for choosing this endpoint, and whether ascertainment was uniform across patients are not described, the outcomes having been ascertained and supplied by the source cohort's authors. \\[2pt]
8b & If outcome assessment requires subjective interpretation, describe the qualifications and demographic characteristics of the outcome assessors & \emph{Not reported} & The endpoint is a documented-rhythm criterion (>=30 s of AF or flutter) ascertained by the source cohort's investigators before this study began, and the manuscript states nothing about who ascertained or adjudicated it, their qualifications, or their demographics. \\[2pt]
8c & Report any actions to blind assessment of the outcome to be predicted & Partially & \textbf{Results Sec. 5.2 sec:clinical} --- Formal blinding of outcome assessors to predictor values is not reported, since outcomes were ascertained independently by the cohort's authors before the predictors existed; what the manuscript does document is the corresponding analyst-side separation, namely that Sec. 8 of the pre-registration was frozen and git-tagged ''before any feature was placed beside that column'' and the analysis was run once and reported as it came out. \\[2pt]
9a & Describe the choice of initial predictors (e.g., literature, previous models, all available predictors) and any pre-selection of predictors before model building & \textbf{Reported} & \textbf{Methods, opening and ''The SFI feature and its competitors''; Results sec:clinical; Sec. 5.4 redundancy table; Sec. 5.5 sec:regions} --- The full 15-column predictor set is specified a priori and enumerated -- six named competitor features (fibrosis burden, fibrosis spatial entropy, fibrosis patch size, Stoer-Wagner global min-cut, percolation threshold, and lambda2-alone, the last included specifically to isolate what SFI adds over bare algebraic connectivity) plus a nine-column SFI block entered strictly as an add-on -- with the baseline set frozen before any label was inspected and no data-driven pre-selection or univariate screening performed (the univariate AUCs of all fifteen are reported only as a post hoc diagnostic); the one predictor-construction choice not pre-registered, the region partition and patient-level reduction, is disclosed as a logged deviation fixed without reference to the outcome. \\[2pt]
9b & Clearly define all predictors, including how and when they were measured (and any actions to blind assessment of predictors for the outcome and other predictors) & Partially & \textbf{Methods, ''The SFI feature and its competitors''; Methods, ''The edge weight''; Methods, ''Substrate''; Methods, ''Cohort flow''; Results, ''Region definition and the patient-level reduction''; Results, clinical endpoint} --- The nine patient-level SFI predictors are defined exactly, together with the edge-weight and fibrosis definitions that generate them and the fact that all predictors are computed from the pre-ablation LGE-MRI-derived mesh, i.e. before ablation and before the two-year outcome window, and outcome-blind (closed-form given the mesh, never fitted to the label, with the protocol frozen and git-tagged ''before any feature was placed beside that column''); however, the six competitor predictors are named only (fibrosis burden, fibrosis spatial entropy, fibrosis patch size, Stoer-Wagner min-cut, percolation threshold, lambda2-alone) without their computational definitions, which are deferred to the released code and pre-registration. \\[2pt]
9c & If predictor measurement requires subjective interpretation, describe the qualifications and demographic characteristics of the predictor assessors & \emph{Not reported} & The predictors as computed in this study are deterministic closed-form functions of the deposited meshes with no human assessor, but the fibrosis input for the clinical cohort is the cohort authors' per-element tissue-tag classification of LGE-MRI (an observer-dependent upstream segmentation), and the manuscript reports nothing about the qualifications, experience, or demographics of the people who performed that segmentation. \\[2pt]
10 & Explain how the study size was arrived at (separately for development and evaluation), and justify that the study size was sufficient to answer the research question. Include details of any sample siz & \textbf{Reported} & \textbf{Methods, ''Cohort flow''; Results, clinical endpoint; Limitations, item 3} --- The clinical study size was not chosen but fixed by the deposit -- all 82 UW patients with supplied outcomes were analysed with no exclusions and no attrition (48 events, 58.5\%) -- and the manuscript explicitly states the size is insufficient, reporting power of 0.24 (0.13 under a weaker shared-noise assumption) at dAUC=0.05, a minimum detectable effect of dAUC approximately 0.10-0.15, and 371-859 patients needed for 80\% power at dAUC=0.05, which is why the 0.05 effect-size gate was pre-specified as not reusable here. \\[2pt]
11 & Describe how missing data were handled. Provide reasons for omitting any data & Partially & \textbf{Methods, ''Cohort flow''; Methods, ''Substrate''; Results, clinical endpoint} --- Outcome data are complete (all 82 patients followed the full two years, no censoring) and no subject is excluded or lost, and the manuscript documents which cohort fields were unavailable and what was substituted (no image-intensity ratios, so fibrosis is derived from tissue tags; no atrial coordinates, so a PCA chart is substituted; no fibre orientations, so the shipped constant vector is used), but there is no explicit statement of whether any predictor values were missing in the analysis matrix and no missing-data or imputation strategy is described. \\[2pt]
12a & Describe how the data were used (e.g., for development and evaluation of model performance) in the analysis, including whether the data were partitioned, considering any sample size requirements & \textbf{Reported} & \textbf{Methods, ''Evaluation pipeline''; Methods, ''Evaluation protocol''; Results, clinical endpoint; Limitations, item 3} --- All 82 patients are used for both development and evaluation through patient-held-out GroupKFold cross-validation (each patient its own group, k=min(5, n\_groups)) with an inner GroupKFold for model selection and out-of-fold probabilities pooled for the DeLong comparison, with no separate held-out or external test set for the primary; the data provider's 67/15 split is reported only as a non-confirmatory analysis and explicitly labelled non-independent and underpowered (8 events), and the sample-size implications of this design are quantified in the power discussion. \\[2pt]
12b & Depending on the type of model, describe how predictors were handled in the analyses (functional form, rescaling, transformation, or any standardisation) & Partially & \textbf{Methods, ''Evaluation pipeline''; Results, ''Reduction to the patient''} --- The manuscript states that the logistic-regression arm is a standardisation-plus-LogisticRegression pipeline (so scaling is fitted within the cross-validation folds) and specifies exactly how the nine per-region SFI values are reduced to patient-level scalars (total, region mean/max/std, top1-top3, edge-fragility max/mean), but it does not state the functional form assumed for continuous predictors (all appear to enter untransformed and linearly, with no discussion of nonlinearity, splines, categorisation, or winsorising) and does not say whether any rescaling is applied in the gradient-boosting arm. \\[2pt]
12c & Specify the type of model, rationale2, all model-building steps, including any hyperparameter tuning, and method for internal validation & \textbf{Reported} & \textbf{Methods, ''Evaluation pipeline''; Methods, ''Evaluation protocol''; Results, clinical endpoint} --- Two model types are specified with rationale -- penalised logistic regression (liblinear, max 2000 iterations) with C selected from \{0.03, 0.1, 0.3, 1.0, 3.0\} in an inner GroupKFold so tuning never sees the outer test fold, and a deliberately untuned gradient-boosted ensemble (120 trees, depth 2, learning rate 0.05, subsample 0.8) justified as a fixed shallow reference because a search at this sample size would fit the folds and logged as a deviation from the pre-registered nested selection for both classifiers -- with internal validation by patient-held-out GroupKFold, at least 10,000 stratified bootstrap resamples, the paired DeLong test (whose approximate validity under pooled out-of-fold scoring is stated), and a Gaussian-noise substitution control; SFI itself is closed-form and contributes no tuning degrees of freedom. \\[2pt]
12d & Describe if and how any heterogeneity in estimates of model parameter values and model performance was handled and quantified across clusters (e.g., hospitals, countries). See TRIPOD-Cluster for addit & \emph{n/a} & \textbf{Results, clinical endpoint; Methods, ''Evaluation pipeline''} --- The clinical endpoint uses a single-source cohort in which each patient contributes one row and is its own cross-validation group, so there are no sites, centres, or repeated-measure clusters across which heterogeneity could be assessed; the manuscript nonetheless characterises heterogeneity across cross-validation folds (event rates 25\% to 81\%, training/test prevalence correlation +0.995) and shows it drives the anti-predictive pooled AUC of 0.250, which within-fold scoring partly removes (0.250 to 0.345). \\[2pt]
12e & Specify all measures and plots used (and their rationale) to evaluate model performance (e.g., discrimination, calibration, clinical utility) and, if relevant, to compare multiple models & Partially & \textbf{Methods, ''Evaluation pipeline'' and ''Evaluation protocol''; Results, Sec. \textbackslash\{\}label\{sec:clinical\}} --- Discrimination and model comparison are fully specified for the clinical endpoint (out-of-fold AUC 0.250 baseline vs 0.447 with SFI, dAUC +0.197 with a >=10,000-resample bootstrap interval [+0.054,+0.334], paired DeLong p=0.0061, plus a parallel gradient-boosting arm and a 20-draw Gaussian-noise control), but no calibration measure or plot (no calibration curve, calibration slope/intercept, or Brier score) and no clinical-utility or decision-curve analysis is reported, and no ROC or calibration figure accompanies the clinical endpoint; calibration is addressed only narratively atwhere miscalibration to training-fold prevalence is diagnosed as the cause of the AUC inversion. \\[2pt]
12f & Describe any model updating (e.g., recalibration) arising from the model evaluation, either overall or for particular sociodemographic groups or settings & \emph{n/a} & \textbf{Results, Sec. \textbackslash\{\}label\{sec:clinical\}} --- No model updating or recalibration was performed: the manuscript explicitly labels the within-fold rescoring that partially removes the prevalence-driven inversion (0.250 -> 0.345) as ''a diagnostic and not a re-analysis'', and the pre-registered fixed-sequence stopping rule halted all further testing after the primary endpoint was not met. \\[2pt]
12g & For model evaluation, describe how the model predictions were calculated (e.g., formula, code, object, application programming interface) & \textbf{Reported} & \textbf{Methods, ''Evaluation pipeline''; Reproducibility} --- Predictions are out-of-fold predicted probabilities pooled across patient-held-out GroupKFold folds, generated by fully specified scikit-learn pipelines (standardisation plus LogisticRegression, liblinear, max 2000 iterations, C selected in an inner GroupKFold over \{0.03,0.1,0.3,1.0,3.0\}; and a frozen untuned gradient-boosted ensemble of 120 trees, depth 2, learning rate 0.05, subsample 0.8), reproducible by running the named script scripts/clinical\_endpoint.py under seeded stochastic steps with reported library versions and commit hash; the specification is via code and hyperparameters rather than published fitted coefficients or an API, since the model is not offered for use. \\[2pt]
13 & If class imbalance methods were used, state why and how this was done, and any subsequent methods to recalibrate the model or the model predictions & Partially & \textbf{Results, Sec. \textbackslash\{\}label\{sec:clinical\}; Methods, ''Evaluation protocol''} --- The outcome distribution is reported (48 events among 82 patients) and severe between-fold event-rate imbalance (25\% to 81\%) is quantified and analysed at length as the mechanism producing the anti-predictive pooled AUC, but the manuscript never states whether any class-imbalance handling (class weighting, resampling, stratified fold construction) was applied in model fitting -- folds are unstratified GroupKFold and the only stratification mentioned is of the bootstrap resamples -- and no recalibration was performed afterwards. \\[2pt]
14 & Describe any approaches that were used to address model fairness and their rationale & \emph{Not reported} & The manuscript contains no fairness assessment and no rationale for its absence: no patient demographics (sex, age, race/ethnicity, AF type, comorbidity) are reported for the 82-patient clinical cohort and no subgroup performance is examined; the only stratified analysis in the clinical section (SFI's failure to separate the 35 fibrillation from the 13 flutter recurrers, is a mechanistic feature check, not a fairness evaluation. \\[2pt]
15 & Specify the output of the prediction model (e.g., probabilities, classification). Provide details and rationale for any classification and how the thresholds were identified & \textbf{Reported} & \textbf{Methods, ''Reporting guideline'' and ''Evaluation pipeline''; Results, Sec. \textbackslash\{\}label\{sec:clinical\}} --- The model output is a continuous predicted probability of two-year recurrence, evaluated threshold-free by AUC and DeLong comparison; no classification threshold is applied and the manuscript states explicitly that it ''reports no sensitivity, specificity, or predictive values at a diagnostic threshold,'' so no threshold rationale is required. \\[2pt]
16 & Identify any differences between the development and evaluation data in healthcare setting, eligibility criteria, outcome, and predictors & \emph{n/a} & \textbf{Methods, ''Evaluation pipeline''; Results, Sec. \textbackslash\{\}label\{sec:clinical\}} --- There is no separate evaluation dataset for the clinical endpoint -- development and evaluation are the same 82-patient UW cohort under patient-held-out cross-validation, so no differences in setting, eligibility, outcome, or predictors can arise -- and the manuscript explicitly flags the one split-based analysis (the data provider's own 67/15 split) as ''explicitly not independent,'' since those 15 patients also enter the primary analysis. \\[2pt]
17 & Name the institutional research board or ethics committee that approved the study and describe the participant-informed consent or the ethics committee waiver of informed consent & \emph{Not reported} & The manuscript contains no ethics statement: it does not name an institutional review board or ethics committee, does not report the human-data exemption determination the authors obtained, and does not state whether informed consent was obtained or waived; it says only that both mesh cohorts are public and de-identified and that the two-year outcomes supplied by the UW cohort's authors carry no confidentiality restriction, and the four mentions of an IRB elsewhere refer to the source cohort's IRB withholding imaging data, not to any approval for this study. \\[2pt]
\midrule
\multicolumn{4}{@{}l}{\textbf{Open science}}\\
\midrule
18a & Give the source of funding and the role of the funders for the present study & \emph{Not reported} & The manuscript contains no funding statement anywhere; the author footnote declares only that this is an ''Independent, unaffiliated in-silico computational-cardiology project,'' which describes affiliation and not the source of funding or any funder role, so the known fact that the study was unfunded must still be added as an explicit statement. \\[2pt]
18b & Declare any conflicts of interest and financial disclosures for all authors & \emph{Not reported} & No conflict-of-interest or financial-disclosure statement appears anywhere in the manuscript (a search of the full 1803-source returns no instance of ''conflict'', ''competing interest'', or ''disclosure''), so the author's declaration of no conflicts must be added explicitly. \\[2pt]
18c & Indicate where the study protocol can be accessed or state that a protocol was not prepared & Partially & \textbf{Methods,; Clinical endpoint,and 688-692; Reproducibility} --- The manuscript states that the full protocol is the pre-registration file \textbackslash\{\}code\{docs/PRE\_REGISTRATION.md\}, that Section 8 of it governing the clinical endpoint was frozen and git-tagged before any feature was placed beside the outcome column, and cites specific protocol clauses (S8.6 baseline-degradation guard, S8.7 stopping rule), but it gives only a repository-relative path with no URL, DOI, or archived location at which a reader or editor could actually access that protocol. \\[2pt]
18d & Provide registration information for the study, including register name and registration number, or state that the study was not registered & Partially & \textbf{Clinical endpoint,; Methods,; Abstract/Keywords,and 115-120, 129} --- The manuscript repeatedly reports that the analysis was pre-registered, frozen and git-tagged before the two-year recurrence outcomes were joined, but it gives no public registry name or registration number (e.g. OSF, ClinicalTrials.gov) and never states whether the study was or was not registered in such a registry, which is what this item requires. \\[2pt]
18e & Provide details of the availability of the study data & \textbf{Reported} & \textbf{Reproducibility,; Methods (Substrate)} --- The Reproducibility section identifies both source datasets as public and de-identified (Roney LA cohort, Zenodo 5801337; UW/Boyle meshes, Dryad 10.5061/dryad.kkwh70sg0), states that the two-year recurrence outcomes used for the clinical endpoint were supplied by the UW cohort's authors who confirm they carry no confidentiality restriction and are available from that group with 15 of the 82 appearing in the source publication, and discloses that the LGE intensity ratios and fibre orientations are withheld under that cohort's IRB and are not obtainable. \\[2pt]
18f & Provide details of the availability of the analytical code7 & Partially & \textbf{Reproducibility,; footnote} --- The manuscript names the exact analysis scripts that reproduce each result (including \textbackslash\{\}code\{scripts/clinical\_endpoint.py\} for the clinical endpoint) and pins the full software stack (Python 3.12.10, NumPy 2.5.1, SciPy 1.18.0, scikit-learn 1.9.0, pandas 3.0.5, Matplotlib 3.11.1, openCARP v19.0), stating that all stochastic steps are seeded and a commit hash accompanies every reported number, but it never gives a URL, DOI, or licence for the ''project repository'' in which that code lives, so the code is not in fact locatable from the paper. \\[2pt]
\midrule
\multicolumn{4}{@{}l}{\textbf{Patient and public involvement}}\\
\midrule
19 & Provide details of any patient and public involvement during the design, conduct, reporting, interpretation, or dissemination of the study or state no involvement & \textbf{Reported} & \textbf{Limitations, item on patient and public involvement} --- The manuscript states there was no patient or public involvement: this is a secondary computational analysis of previously published de-identified deposits with no participant contact at any stage, and the endpoint definition was set by the originating study rather than with patient input. TRIPOD+AI item 19 is satisfied by an explicit statement of no involvement. \\[2pt]
\midrule
\multicolumn{4}{@{}l}{\textbf{Results}}\\
\midrule
20a & Describe the flow of participants through the study, including the number of participants with and without the outcome and, if applicable, a summary of the follow-up time. A diagram may be helpful 7 T & Partially & \textbf{Clinical endpoint,; Methods (Cohort flow)} --- For the clinical endpoint the manuscript states that per-patient two-year post-ablation outcomes were supplied for all 82 UW patients, defines recurrence as at least 30 s of documented atrial fibrillation or flutter after a 90-day blanking period, states that every patient was followed the full two years, gives 48 events at n=82 with the 48 recurrers split 35 fibrillation / 13 flutter, and the Methods cohort-flow paragraph explains that the 164 released UW meshes are 82 patients imaged twice with only the 82 pre-ablation meshes analysed and ''no attrition to report''; however, the 34 non-recurrers and the 58.5\% event rate are never stated explicitly, loss to follow-up/censoring is never named as such, and no participant flow diagram is provided. \\[2pt]
20b & Report the characteristics overall and, where applicable, for each data source or setting, including the key dates, key predictors (including demographics), treatments received, sample size, number of & Partially & \textbf{Methods, ''Cohort flow'' and ''Substrate'' paragraphs; Results, sec:clinical; Limitations item 3} --- For the clinical endpoint the manuscript reports sample size (82 UW patients, all pre-ablation meshes, no exclusions or attrition), outcome events (48 recurrences, of which 35 atrial fibrillation and 13 flutter), complete 2-year follow-up with no censoring, the treatment context (all patients ablated, with pre- and post-ablation imaging), the predictor set and its availability limits (UW ships no image-intensity ratios, fibre orientations or atrial coordinates), and the data source, but it gives no demographic or clinical characteristics (age, sex, race/ethnicity, AF type, comorbidities, LA size) and no key dates of imaging, ablation or follow-up, and provides no per-data-source characteristics table. \\[2pt]
20c & For model evaluation, show a comparison with the development data of the distribution of important predictors (demographics, predictors, and outcome). & \emph{n/a} & \textbf{Results, sec:clinical; Methods, ''Evaluation pipeline''} --- The clinical model was developed and evaluated within the same 82-patient cohort by patient-held-out cross-validation with no separate development dataset and no external evaluation cohort, so a development-versus-evaluation comparison of predictor and outcome distributions does not arise; the manuscript does note that the one available non-independent split (the data provider's own 67/15 partition) is explicitly not independent because those 15 patients also enter the primary analysis. \\[2pt]
21 & Specify the number of participants and outcome events in each analysis (e.g., for model development, hyperparameter tuning, model evaluation) & \textbf{Reported} & \textbf{Results, sec:clinical; Limitations item 3} --- Every clinical analysis states its denominator and event count: the primary analysis uses all 82 patients with 48 events, the per-fold event rates are given (25\% to 81\%), the univariate screen covers all fifteen features on the same 82, the Gaussian-noise control reports 20 draws with 6 of 20 matching or exceeding the observed increment, and the two non-confirmatory analyses state their own counts (the 67/15 split with 8 events; the 35-versus-13 fibrillation/flutter comparison among the 48 recurrers). \\[2pt]
22 & Provide details of the full prediction model (e.g., formula, code, object, API) to allow predictions in new individuals and to enable third-party evaluation and implementation, including any restricti & Partially & \textbf{Methods, ''Evaluation pipeline'' and ''The SFI feature and its competitors''; Results, sec:regions; Reproducibility} --- The manuscript fully specifies the feature definitions (six competitors plus the nine-column SFI block, each defined in closed form from the mesh), the modelling pipeline and all hyperparameters (standardisation plus liblinear logistic regression with C selected in an inner GroupKFold over a stated grid; frozen gradient boosting at 120 trees, depth 2, learning rate 0.05, subsample 0.8), seeds, software versions and the executable script (scripts/clinical\_endpoint.py), and it discloses the access restrictions (LGE intensity ratios and fibre orientations are IRB-withheld and unobtainable; the recurrence outcomes are available from the cohort authors and carry no confidentiality restriction), but it presents no fitted coefficients, intercept or final locked model, because the endpoint was not met and no model is put forward for use in new individuals. \\[2pt]
23a & Report model performance estimates with confidence intervals, including for any key subgroups (e.g., sociodemographic). Consider plots to aid presentation & Partially & \textbf{Results, sec:clinical} --- The clinical results are given as out-of-fold AUCs with a 95\% bootstrap confidence interval and a p-value on the primary increment (baseline 0.250, plus-SFI 0.447, dAUC +0.197 [+0.054, +0.334], DeLong p=0.0061), plus the gradient-boosting arm (baseline 0.516, dAUC +0.021, p=0.738, no interval given), the univariate AUC range across all fifteen features (0.435-0.547), and two subgroup/split analyses (67/15 split: +0.196 and -0.107; fibrillation versus flutter among recurrers: -0.009 and -0.044), but no confidence intervals accompany the individual AUCs or the subgroup estimates, and no calibration measures, ROC curve or any other plot is presented for the clinical endpoint (all figures in the paper are in-silico). \\[2pt]
23b & If examined, report results of any heterogeneity in model performance across clusters. See TRIPOD Cluster for additional details & Partially & \textbf{Results, sec:clinical} --- The only clustering examined is the cross-validation fold structure, and its heterogeneity is reported in detail (per-fold event rates from 25\% to 81\%, a +0.995 correlation between training- and test-fold prevalence, and a within-fold-scored baseline of 0.345 against the pooled 0.250, reported explicitly as a diagnostic); the clinical cohort is from a single centre and no performance heterogeneity across sites, demographic groups or other clinically meaningful clusters was examined or reported. \\[2pt]
24 & Report the results from any model updating, including the updated model and subsequent performance & \emph{n/a} & \textbf{Results, sec:clinical} --- No model updating, recalibration or revision was performed: the within-fold rescoring that partially reverses the pooled-AUC inversion is labelled explicitly as ''a diagnostic and not a re-analysis,'' and the pre-registered fixed-sequence stopping rule halted testing at the first non-rejection so that no further model was fitted or updated. \\[2pt]
\midrule
\multicolumn{4}{@{}l}{\textbf{Discussion}}\\
\midrule
25 & Give an overall interpretation of the main results, including issues of fairness in the context of the objectives and previous studies & Partially & \textbf{Discussion; Results, sec:clinical; Conclusion} --- The manuscript interprets the clinical result against its pre-registered objective without overstatement, concluding that the primary endpoint is not met and that the arithmetically significant +0.197 increment is an artefact of an anti-predictive baseline that nine columns of Gaussian noise reproduce, and it explicitly bounds what is and is not claimed; comparison with previous work is limited to the pre-committed rule that the baseline be judged against the source study's published performance without quoting that study's numbers, and no fairness, equity or demographic-subgroup assessment of the model is reported anywhere. \\[2pt]
26 & Discuss any limitations of the study (such as a non-representative sample, sample size, overfitting, missing data) and their effects on any biases, statistical uncertainty, and generalizability & \textbf{Reported} & \textbf{Limitations items 3, 4 and 5; Results, sec:clinical; Discussion} --- The limitations of the clinical analysis are stated directly: the cohort is underpowered by construction (power 0.24 at best at the simulator gate, minimum detectable effect dAUC 0.10-0.15, 371-859 patients needed for 80\% power at 0.05), the sample is non-representative for this question because UW fibrosis is categorical rather than continuous and so degrades the very baseline SFI must beat, the observed increment is attributed to an over-fitted anti-predictive baseline diluted by any nine columns, and the availability limits of the input data (withheld intensity ratios, placeholder fibre field, PCA surrogate for atrial coordinates) are documented with their effects; there were no missing outcomes, since all 82 patients were followed the full two years. \\[2pt]
\midrule
\multicolumn{4}{@{}l}{\textbf{Other information}}\\
\midrule
27a & Describe how poor quality or unavailable input data (e.g., predictor values) should be assessed and handled when implementing the prediction model & \emph{n/a} & No prediction model is proposed for clinical implementation, since the clinical endpoint was not met and its apparent increment is reported as an artefact, so guidance on handling poor-quality or unavailable inputs at the point of use does not arise; for completeness, the substitutions made for inputs the release withholds (tag-derived fibrosis in place of image-intensity ratios, a PCA chart in place of universal atrial coordinates) are documented in Methods. \\[2pt]
27b & Discuss whether users will be required to interact in the handling of the input data or use of the model, and what level of expertise is required of users & \emph{n/a} & No model is offered for use, so the manuscript specifies no user interaction with input data or with the model and states no required level of user expertise; the only usage instructions given are for reproducing the reported analysis from the released scripts (Reproducibility,. \\[2pt]
27c & Discuss any next steps for future research, with a specific view to applicability and generalizability of the model & \textbf{Reported} & \textbf{Limitations item 3; Discussion; Conclusion} --- The manuscript states what a next study would require -- a cohort with continuous fibrosis imaging and several hundred patients (371-859 for 80\% power at the 0.05 gate) -- explicitly declines to read the clinical null as evidence that SFI adds nothing to recurrence prediction, notes that a small effect remains compatible with the data, and generalises the portable part of the work by giving the label-free validity screen as an a priori check for any Fiedler-based marker in other domains while cautioning that it bounds estimator accuracy rather than predictive content. \\[2pt]
\bottomrule
\end{longtable}

\end{document}
